\section{Contoh Penilaian Proyek Akhir}
\label{sec:contoh-bab17}

Agar rubrik pada Tabel~\ref{tab:rubrik-proyek} tidak menjadi
dokumen yang abstrak, berikut contoh penerapannya pada satu proyek.

\subsection{Deskripsi Proyek yang Dinilai}

Seorang mahasiswa mengerjakan Pilihan 1, yaitu kalkulator dengan
validasi masukan. Proyeknya memuat:

\begin{itemize}
    \item prosedur \code{BacaAngka} yang membaca string masukan
          dengan \code{INT 21h} fungsi \hex{0A} lalu mengubahnya
          menjadi nilai biner;
    \item prosedur \code{Hitung} yang menerima dua parameter melalui
          stack beserta kode operator, lalu mengembalikan hasil pada
          \reg{AX} dan \reg{DX};
    \item prosedur \code{TulisAngka} yang mengubah hasil biner
          kembali menjadi deretan digit dan mencetaknya;
    \item penanganan pembagian dengan nol yang mencetak pesan
          kesalahan dan kembali ke menu.
\end{itemize}

\subsection{Penerapan Rubrik}

Tabel~\ref{tab:contoh-penilaian} memperlihatkan cara pemberian skor
beserta alasan yang harus dituliskan penilai. Kolom alasan inilah
yang membuat penilaian dapat dipertanggungjawabkan.

\begin{table}[H]
    \centering
    \caption{Contoh penerapan rubrik proyek akhir.}
    \label{tab:contoh-penilaian}
    \footnotesize
    \begin{tabular}{@{}p{2.9cm}cp{7.6cm}@{}}
        \toprule
        \textbf{Kriteria} & \textbf{Skor} & \textbf{Alasan} \\
        \midrule
        Kebenaran fungsi & 4 &
        Empat operasi berjalan benar; pembagian dengan nol ditangani
        dengan pesan kesalahan; masukan bukan angka ditolak \\
        \addlinespace
        Kualitas kode & 3 &
        Prosedur terpisah dengan baik, tetapi prosedur
        \code{TulisAngka} memiliki tiga baris yang berulang dan tidak
        dijadikan subprosedur \\
        \addlinespace
        Penggunaan instruksi & 4 &
        Perkalian pangkat dua memakai \asm{SHL} alih-alih
        \asm{MUL}; mode \eng{base-relative} dipakai untuk menelusuri
        array digit \\
        \addlinespace
        Pengujian dan \eng{debugging} & 3 &
        Tabel \eng{trace} disertakan untuk dua dari tiga bug; satu
        bug terakhir ditemukan dengan bantuan dosen \\
        \addlinespace
        Laporan dan presentasi & 4 &
        Laporan memuat rancangan, kendala, dan bukti pengujian;
        presentasi menjelaskan alasan pemilihan mode pengalamatan \\
        \midrule
        \textbf{Rata-rata} & \textbf{3{,}6} &
        Setara dengan kategori ``baik'' menuju ``sangat baik'' \\
        \bottomrule
    \end{tabular}
\end{table}

\subsection{Pelajaran dari Contoh}

Tiga hal dapat dipelajari dari contoh di atas.

\begin{enumerate}
    \item \textbf{Skor tertinggi tidak menuntut proyek yang paling
          rumit.} Proyek sederhana yang dikerjakan dengan benar dan
          didokumentasikan dengan baik memperoleh nilai lebih tinggi
          daripada proyek rumit yang tidak dapat dijelaskan.

    \item \textbf{Alasan harus dituliskan.} Rubrik tanpa alasan
          hanya menghasilkan angka; rubrik dengan alasan
          menghasilkan umpan balik yang dapat ditindaklanjuti.

    \item \textbf{Bukti menentukan.} Kriteria pengujian dan
          \eng{debugging} tidak dapat dinilai tanpa tabel
          \eng{trace} dan catatan bug. Inilah alasan jurnal
          praktikum wajib dipelihara sejak pertemuan pertama.
\end{enumerate}

\begin{infobox}
Perhatikan bahwa seluruh kriteria pada rubrik berasal dari CPMK.
Kriteria ``penggunaan instruksi'' misalnya berasal dari \CPMK{02},
sedangkan ``kebenaran fungsi'' dan ``kualitas kode'' berasal dari
\CPMK{03}. Tidak ada kriteria yang tidak dapat ditelusuri ke CPMK
--- inilah yang dimaksud dengan asesmen berbasis OBE.
\end{infobox}
